\bookchapter{الفصل 43: هندسة الحوار — تقنيات الحوار للتعاون الفعال مع الذكاء الاصطناعي}

\chapterintro{الحوار هو طريقة التفاعل الأساسية في البرمجة بالذكاء الاصطناعي. لكن انتبه: “الحوار” لا يعني “الدردشة” — فالحوار الفعال في البرمجة بالذكاء الاصطناعي تقنية تواصل منظمة نسميها “هندسة الحوار”. وهي تدرس: كيف تحصل على أجود مخرجات برمجية بأقل عدد من جولات الحوار. بصراحة، هذا الفصل يعلّمك كيف “تقول كلامك بوضوح” للذكاء الاصطناعي.}

\booksection{أولاً: استراتيجية طبقات الموجه (Prompt)}

يمكن تقسيم الحوار مع الذكاء الاصطناعي إلى ثلاث طبقات، لكل طبقة وظيفة مختلفة. حين تفهم هذه الطبقات الثلاث ستتمكن من تركيب موجهات فعالة كما تركب المكعبات.

\begin{center}
\begin{tabularx}{\linewidth}{|X|X|X|X|}
\hline
\textbf{الطبقة} & \textbf{الوظيفة} & \textbf{المثال} & \textbf{التكرار} \tabularnewline
\hline
طبقة النظام & تحديد دور الذكاء الاصطناعي والقيود الشاملة & "أنت خبير React، ويجب كتابة الكود بـ TypeScript" & تضبط مرة واحدة عند بدء المشروع \tabularnewline
طبقة المهمة & وصف المهمة المحددة المطلوب إنجازها & "نفذ مكوّن نموذج يتضمن حقلَي اسم المستخدم وكلمة المرور" & يقدم عند بدء كل مهمة \tabularnewline
طبقة الملاحظات & تقييم مخرجات الذكاء الاصطناعي وتصحيحها & "حقل كلمة المرور ينقصه زر إظهار/إخفاء، أضفه من فضلك" & يقدم حسب جودة المخرجات \tabularnewline
\hline
\end{tabularx}
\end{center}

\begin{center}\small\color{gray}الطبقات الثلاث لموجه الحوار\end{center}

تكتب موجهات طبقة النظام عادة في ملف ذاكرة المشروع (مثل CLAUDE.md) ولا يلزم تكرارها في كل حوار. أما طبقة المهمة فهي الأكثر استخداما في الحوار — ويجب أن تتضمن ثلاثة عناصر: “ماذا نفعل” و“بأي تقنية” و“ما القيود”. وطبقة الملاحظات هي الأكثر اختبارا لمهارتك — فالملاحظات الجيدة تجعل الذكاء الاصطناعي يصحح بدقة، والملاحظات السيئة تجعله يعيد التخمين فقط.

\begin{tcbinfo}{نصيحة: صيغة الملاحظات الثلاثية}
قالب عملي للملاحظات هو “حدد المشكلة + اشرح المتوقع + قدم السياق”. مثلا: “معالجة الخطأ في السطر 25 تستخدم console.log فقط (تحديد المشكلة)، وبيئة الإنتاج ينبغي أن تستخدم معالج أخطاء موحدا وترجع رسالة مفهومة للمستخدم (شرح المتوقع)، ويوجد في المشروع دالة handleError جاهزة للاستخدام (تقديم السياق)”. هذه الملاحظات أنجح عشر مرات من “عدل معالجة الأخطاء”.

\end{tcbinfo}

\booksection{ثانياً: حوار التصحيح — كيف تصف الخطأ ليحدده الذكاء الاصطناعي بسرعة}

التصحيح (Debugging) من أكثر السيناريوهات تكرارا في البرمجة بالذكاء الاصطناعي. وحين تصطدم بخطأ برمجي فإن طريقة وصفك للمشكلة للذكاء الاصطناعي تحدد مباشرة ما إذا كان سيتمكن من تحديده وإصلاحه بسرعة. إليك النقاط الأربع لحوار التصحيح:

\begin{enumerate}
  \item أعط الذكاء الاصطناعي المعلومات الأصلية "ليراها" بنفسه، لا روايتك عنها. انقل إليه سجل الأخطاء وحزمة الاستثناءات والكود ذا الصلة مباشرة بالإشارة إليه (@)، بدلا من وصفه بعبارة "يبدو أن هناك خطأ ما". الذكاء الاصطناعي يفهم المعلومات الأصلية بدقة أكبر من فهمه لتفسيرك.
  \item صف الفرق بين "السلوك المتوقع" و"السلوك الفعلي". "أتوقع أن ينتقل التطبيق إلى الصفحة الرئيسية بعد الضغط على الزر (متوقع)، لكن لا يحدث أي رد فعل بعد الضغط فعليا (فعلي)" — هذا الوصف أنفع بكثير من "الزر معطل".
  \item قدم خطوات إعادة الإنتاج. "أدخل اسم المستخدم أولا، ثم اضغط زر تسجيل الدخول، فيظهر هذا الخطأ" — خطوات إعادة الإنتاج تمكن الذكاء الاصطناعي من محاكاة مسار عملك، وتسهل تحديد المشكلة.
  \item وضح ما جربته من قبل. "تحققت من طلب الشبكة، وأعاد 200، والبيانات سليمة" — هذا يساعد الذكاء الاصطناعي على استبعاد المسارات المؤكدة، وألا يهدر وقته في تكرار الفحص.
\end{enumerate}

وصف تصحيح رديء

\begin{lstlisting}
الزر لا يستجيب عند الضغط، افحصه لي
\end{lstlisting}

وصف تصحيح جيد

\begin{lstlisting}
الصفحة لا تستجيب بعد الضغط على زر تسجيل الدخول. السلوك المتوقع: الانتقال إلى صفحة /home بعد الضغط السلوك الفعلي: لا يحدث أي استجابة بعد الضغط، ولا توجد أخطاء في وحدة التحكم خطوات إعادة الإنتاج: 1. أدخل اسم المستخدم admin وكلمة المرور 123456 في صفحة تسجيل الدخول 2. اضغط زر "تسجيل الدخول" 3. الصفحة لا تستجيب ما تم فحصه: - طلب الشبكة طبيعي ويعيد 200 - حدث onClick مرتبط - فحصت كود مكوّن LoginButton الملفات ذات الصلة: @src/components/LoginButton.tsx @src/pages/Login.tsx
\end{lstlisting}

\booksection{ثالثاً: استراتيجية التكرار — خطوات صغيرة سريعة مقابل إعادة هيكلة كبيرة}

هناك استراتيجيتان أساسيتان للتكرار في البرمجة بالذكاء الاصطناعي، تصلحان لسيناريوهات مختلفة. فهم الفرق بينهما يساعدك على الاختيار الصائب في كل حالة.

تصلح استراتيجية “الخطوات الصغيرة السريعة” لتطوير الميزات وإصلاح الأخطاء. جوهرها: اطلب من الذكاء الاصطناعي تعديلا صغيرا واحدا في كل مرة، وبعد نجاح التحقق انتقل إلى التعديل التالي. ميزتها أن كل خطوة قابلة للتحقق والتراجع فتبقى المخاطر تحت السيطرة؛ وعيبها أنها أبطأ وتتطلب تفاعلا حواريا متكررا. تصلح لـ: تطوير ميزات جديدة، والاستكشاف حين لا تكون متأكدا من الحل الأمثل، والتكرار الحذر حين لا تطمئن كثيرا إلى مخرجات الذكاء الاصطناعي.

تصلح استراتيجية “إعادة الهيكلة الكبيرة” للتعديلات الواسعة حين يكون الهدف واضحا. جوهرها: وصف الهدف الكامل دفعة واحدة وترك الذكاء الاصطناعي يخطط التنفيذ بأكمله وينفذه ذاتيا. ميزتها السرعة، إذ يرى الذكاء الاصطناعي الصورة الكلية فيتخذ قرارات معمارية أفضل؛ وعيبها ارتفاع المخاطر، فإذا أخطأ الذكاء الاصطناعي في الفهم كانت تكلفة التصحيح كبيرة. تصلح لـ: ترقيات معمارية واضحة، وترحيل كود واسع النطاق، وإعادة هيكلة مدعومة باختبارات كافية.

\begin{center}
\begin{tabularx}{\linewidth}{|X|X|X|}
\hline
\textbf{البعد} & \textbf{خطوات صغيرة سريعة} & \textbf{إعادة هيكلة كبيرة} \tabularnewline
\hline
حجم التعديل في كل مرة & ميزة واحدة أو خطأ برمجي واحد & وحدة كاملة أو طبقة معمارية كاملة \tabularnewline
تكرار التحقق & التحقق فورا بعد كل تعديل & تحقق موحد بعد اكتمال جميع التعديلات \tabularnewline
جولات الحوار & كثيرة (5-10 جولات لكل ميزة) & قليلة (1-3 جولات لكل وحدة) \tabularnewline
المخاطر & منخفضة، وكل خطوة قابلة للتراجع & مرتفعة، وتكلفة تصحيح الانحراف كبيرة \tabularnewline
السيناريوهات المناسبة & الميزات الجديدة والتطوير الاستكشافي & إعادة الهيكلة والترحيل بهدف واضح \tabularnewline
الشروط المسبقة & لا متطلبات خاصة & يتطلب تغطية اختبارات \tabularnewline
\hline
\end{tabularx}
\end{center}

\begin{center}\small\color{gray}مقارنة بين استراتيجيتي التكرار\end{center}

عمليا تستخدم الاستراتيجيتان معا غالبا. مثلا عند إعادة هيكلة وحدة ما، ابدأ بـ“إعادة الهيكلة الكبيرة” ليقترح الذكاء الاصطناعي الخطة الشاملة، وبعد مراجعتها نفذها تدريجيا بأسلوب “الخطوات الصغيرة السريعة”. هكذا تحافظ على الرؤية الكلية وتضبط مخاطر التنفيذ.

\booksection{رابعاً: إدارة السياق في الحوارات متعددة الجولات}

مع ازدياد جولات الحوار عليك أن تدير سياق الحوار بوعي لتجنب مشكلة “نسيان الذكاء الاصطناعي للمعلومات السابقة”. إليك عدة حيل عملية لإدارة السياق.

أولا، “ثبت” المعلومات الأساسية. اكتب القيود الأساسية بوضوح في بداية الحوار، وأعد التذكير بها في الوقت المناسب خلال الحوار اللاحق. مثلا “تذكر أن بياناتنا مخزنة في localStorage ولا يمكننا استخدام قاعدة بيانات” — أعد ذكر هذا القيد في طلبات التعديل ذات الصلة لمنع الذكاء الاصطناعي من نسيانه.

ثانيا، استخدم بحكمة أسلوب “الملخص ثم المتابعة”. حين يطول الحوار اطلب من الذكاء الاصطناعي “تلخيص ما وصلنا إليه حتى الآن والمهام المتبقية”، ثم افتح حوارا جديدا واجعل هذا الملخص نقطة البداية. هكذا تحتفظ بالسياق الأساسي وتتجنب تدهور الحوارات الطويلة.

ثالثا، “حوار واحد لمهمة واحدة”. هذه أبسط استراتيجيات إدارة السياق وأكثرها فعالية. اجعل كل حوار يعالج مهمة واحدة أو ميزة واحدة، وبعد الانتهاء افتح حوارا جديدا للتالية. هذا يحد طول الحوار طبيعيا ويحافظ على وضوح السياق.

رابعا، استفد من ميزات إدارة السياق في الأدوات. تضم معظم أدوات البرمجة بالذكاء الاصطناعي زر “مسح السياق” أو “جلسة جديدة” — وحين تشعر بتدهور جودة إجابات الذكاء الاصطناعي لا تتردد: افتح جلسة جديدة. فقدان جزء من سجل الحوار أفضل بكثير من مواصلة الصراع داخل سياق متدهور.

\begin{tcbpractice}{لنطبق عملياً: كن "المريض" أمام ذكائك الاصطناعي مرة}
اصنع عمدا صفحة صغيرة في المتصفح وخلق فيها خطأ برمجيا (مثل زر لا يستجيب). ثم صف للذكاء الاصطناعي وصفا كاملا وفق صيغة “وصف التصحيح الجيد” في هذا الفصل: المعلومات الأصلية، المتوقع مقابل الفعلي، خطوات إعادة الإنتاج، وما تم فحصه. اختبر بنفسك فرق السرعة في تحديد المشكلة — وستفهم لماذا يعد “وصف المشكلة” حرفة بحد ذاتها.

\end{tcbpractice}

\begin{tcbsummary}{خلاصة الفصل}
\begin{itemize}
  \item طبقات الموجه الثلاث: طبقة النظام (تكتب في ملف الذاكرة)، طبقة المهمة (ماذا نفعل + التقنية + القيود)، طبقة الملاحظات (تحديد المشكلة + المتوقع + السياق)
  \item نقاط التصحيح الأربع: تقديم المعلومات الأصلية، المتوقع مقابل الفعلي، خطوات إعادة الإنتاج، ما تم فحصه
  \item الخطوات الصغيرة السريعة: مخاطر منخفضة وجولات كثيرة، مناسبة للميزات الجديدة؛ إعادة الهيكلة الكبيرة: مخاطر مرتفعة وجولات قليلة، وتحتاج إلى تغطية اختبارات
  \item حيل إدارة السياق الأربع: تثبيت المعلومات الأساسية، الملخص ثم المتابعة، حوار واحد لمهمة واحدة، وفتح جلسة جديدة بحزم
\end{itemize}

\end{tcbsummary}
